\documentclass{rapportCS} % Ensure rapportCS.cls is available
\usepackage{lipsum}
\usepackage{svg}
\usepackage{float} % Added for 'H' float specifier
\title{Report - Template}
\begin{document}

%--------------------

\logoentreprise{logos/logo.png}

\titre{Transportation's Systems 20-5751}

\mention{Prof. Zahra Amini} 
\trigrammemention{Project}
\master{Keivan Jamali}

\eleve{Project Part 3: One-Dimensional Optimization Methods}

\dates{1405/02/08}


%------------------------------
        
\fairemarges 
\fairepagedegarde 

%----------- Abstract -------------------
%\vspace*{\stretch{1}}
%\begin{center}
	%\begin{abstract}
        %\lipsum[1]
        %\rule{\linewidth}{0.2 mm} \\[0.4 cm]
        %\begin{center}\textbf{Summary :}\end{center} 
        %\lipsum[2]
    %\end{abstract}
%\end{center}
%\vspace*{\stretch{1}}
%\newpage

%----------------------------

%====================================================
\section{Introduction}

In this part, you will implement two fundamental one-dimensional optimization algorithms and apply them to transportation-related problems based on the Sioux Falls network.

The goal is to prepare you for more advanced algorithms (e.g., Frank–Wolfe) by first mastering line search and root-finding techniques.

%------------------------------------------------------
\subsection*{Problem 1: Bisection Method for Flow Equilibrium on a Link}

In transportation systems, congestion effects are often modeled using nonlinear travel time functions. Consider the BPR function:

$$ T(f) = T^0 \left(1 + B \left(\frac{f}{C}\right)^4 \right) $$

We define the total cost on a link as:

$$ Z(f) = f \cdot T(f) $$

Your task is to find the flow value \( f^* \) such that:

$$ \frac{dZ(f)}{df} = K $$

Where:
\begin{itemize}
    \item \( K \) is a constant derived from your student ID
    \item \( C = 1000 \)
    \item \( T^0 \) and \( B \) are generated using your student ID
\end{itemize}

\textbf{Instructions:}
\begin{enumerate}
    \item Derive \( \frac{dZ(f)}{df} \)
    \item Reformulate the problem as a root-finding problem:
    $$ g(f) = \frac{dZ(f)}{df} - K = 0 $$
    \item Use the \textbf{Bisection Method} to solve for \( f^* \)
    \item Use an interval \( [0, C] \)
    \item Stop when:
    $$ |g(f)| < 10^{-5} $$
\end{enumerate}

\textbf{Student-Specific Parameters:}

Let your student ID be \( ID \):
\begin{itemize}
    \item \( seed = int(\text{hashlib}.sha256(\text{student\_id}.encode()).hexdigest(), 16) \% 1e8 \)
    \item \( B = 0.15 + \frac{(int(seed) \% 10)}{100} \)
    \item \( T^0 = 5 + (int(seed) \% 10) \)
    \item \( K = T^0 + (int(seed) \% (T^0*(1+5*B) - T^0)) \)
\end{itemize}

\subsubsection*{IMPORTANT NUMBER}

\(x^* = round(f^*)\) and \(y^* = round(Z(f^*))\). Compute:
$$ R_1 = int(hash(tuple([x^*, y^*])) \% 100000) $$

%------------------------------------------------------
\subsection*{Problem 2: Golden Section Search for Optimal Step Size}

This problem mimics the line search step used in traffic assignment algorithms such as Frank–Wolfe.

You are given a subset of links from the Sioux Falls network. For each link \( (i,j) \), two flow values are provided:
\begin{itemize}
    \item Current flow \( x_{ij} \)
    \item Auxiliary flow \( y_{ij} \)
\end{itemize}

Define the flow on each link as a function of \( \lambda \):

$$ f_{ij}(\lambda) = x_{ij} + \lambda (y_{ij} - x_{ij}), \quad \lambda \in [0,1] $$

The objective function is the total system travel time:

$$ Z(\lambda) = \sum_{(i,j)} f_{ij}(\lambda) \cdot T_{ij}(f_{ij}(\lambda)) $$

where the travel time function follows the BPR form:

$$ T_{ij}(f) = T_{ij}^0 \left( 1 + B_{ij} \left( \frac{f}{C} \right)^4 \right), \quad C = 1000 $$

\textbf{Input Data:}

For each student, a data file is provided. This file contains:
\begin{itemize}
    \item A list of selected links \( (i,j) \)
    \item For each link:
    \begin{itemize}
        \item \( x_{ij} \): current flow
        \item \( y_{ij} \): auxiliary flow
        \item \( T_{ij}^0 \): free-flow travel time
        \item \( B_{ij} \): BPR parameter
    \end{itemize}
\end{itemize}

Each student must use the dataset corresponding to their \texttt{student\_id}.

\textbf{Instructions:}
\begin{enumerate}
    \item Read the provided data file corresponding to your \texttt{student\_id}
    \item Implement the \textbf{Golden Section Search} algorithm
    \item Minimize \( Z(\lambda) \) over \( \lambda \in [0,1] \)
    \item Stop when:
    $$ |b - a| < 10^{-5} $$
\end{enumerate}

\subsubsection*{IMPORTANT NUMBER}

Let \( \lambda^* \) be the optimal value rounded to 6 decimal places and \( y^* = round(Z(\lambda^*)) \). Compute:

$$ R_2 = int(hash(tuple([x_golde, y_golde])) \% 100000) $$
%------------------------------------------------------
\subsection*{Submission for Part 3}

Submit:
\begin{itemize}
    \item Your implementation of:
    \begin{itemize}
        \item Bisection Method.
        \item Golden Section Search.
        \item The code to solve problem 1 and problem 2 using the above methods.
    \end{itemize}
    \item A short report including:
    \begin{itemize}
        \item Derived formulas
        \item Final values of \( f^* \) and \( \lambda^* \)
    \end{itemize}
    \item A single text file containing IMPORTANT NUMBERS:
    \begin{itemize}
        \item \( R_1 \)
        \item \( R_2 \)
    \end{itemize}
\end{itemize}    
%------------------------------------------------------
%\bibliographystyle{plain}
%\bibliography{references} 

\end{document}
